\documentclass[10pt,a4paper]{amsart}
\usepackage[margin=19mm]{geometry}
\usepackage{amsmath,amssymb,mathtools}
\usepackage[scr=rsfs,cal=euler]{mathalfa}
\usepackage{xfrac}
\usepackage[unicode,hidelinks,bookmarks=false]{hyperref}
\newcommand{\ie}{i.e.\ }
\newcommand{\cf}{cf.\ }
\newcommand{\resp}{respectively\ }
\newcommand{\Subsection}{\subsection}
\providecommand{\nobreakdash}{\nobreak\hbox{-}}
\setlength{\emergencystretch}{2em}
\multlinegap0pt
\usepackage{tikz}
\newtheorem{lemma}{Lemma}
\title{Point evaluation for polynomials on the circle}
\author{Sarah May Instanes}
\date{Source: arXiv:2509.22035v2 (CC BY 4.0)}
\begin{document}
\maketitle

\noindent\emph{Source and licence.} Point evaluation for polynomials on the circle. Version arXiv:2509.22035v2 (CC BY 4.0), distributed by arXiv under the Creative Commons Attribution 4.0 International licence, http://creativecommons.org/licenses/by/4.0/, as recorded in the arXiv licence metadata for this exact version and shown on the abstract page https://arxiv.org/abs/2509.22035v2. Full licence notice: \texttt{licenses/arxiv-2509.22035-CC-BY-4.0.txt}.\par\smallskip

\noindent\textit{Editorial scope.} Counted unit: the article's lemma lem:structure (source0.tex lines 1182--1188). The article's complete original proof of it is delivered with it (source0.tex lines 1190--1216, one \texttt{proof} environment of 27 lines). No mathematics is written, repaired or re-derived: the statement and the proof are the article's own lines, in the article's own notation, and no retained line is substituted or re-flowed.  Retained uncounted as the source-local context the proof needs: lines 80--82.  Nothing was omitted from any retained span, so the package carries no omission record.  No retained line contains a citation, so the wrapper carries no bibliography and no bibliography entry.  No retained line contains a figure environment, an image-inclusion command or drawing code, so no image is delivered and the package declares no asset dependency.  Editorial formatting only, in addition: an AMS \texttt{amsart} preamble that restates the article's own declarations and macros used by the retained lines, namely the environments \texttt{lemma} on separate counters, together with the standard packages those lines need.  The retained environments print in this document as Lemma 1 in order of appearance, whereas the article prints the same items with the numbering of its own full document.  The complete native source of the article as distributed in the arXiv e-print for this exact version is bundled unchanged as \texttt{sources/185785/source0.tex}.
	\begin{equation}\label{eq:extprob}
		\frac{1}{\mathscr{C}_{d,p}} = \inf_{P \in \mathscr{P}_d} \{ \|P\|_p^p : P(1)=1\},
	\end{equation}

	\begin{lemma}\label{lem:structure}
		Let $1<p<\infty$. Let $d \geq 1$, and let $k=\lfloor d/2 \rfloor$. Then the unique solution of \eqref{eq:extprob} is of the form
		\[\varphi_{d, p}(z)= C\prod_{j=1}^{k} (z-e^{i t_j})(z-e^{-i t_j}),\]
		for $d$ even and 
		\[\varphi_{d, p}(z)=C(z+1) \prod_{j=1}^{k} (z-e^{i t_j})(z-e^{-i t_j}),\]
		for $d$ odd, where $0 < t_1 \leq t_2 \leq \dots \leq t_{k} \leq \pi$ and $C$ is a normalization constant ensuring that $\varphi_{d,p}(1)=1$.
	\end{lemma}

	\begin{proof}
		Assume $P$ is a solution of \eqref{eq:extprob} of degree $n<d$ and let $Q(z)=zP(z)$ be a polynomial of degree at most $d$. Then $\|Q\|_p=\|P\|_p$ and $Q(1)=P(1)=1$. However, since the extremal function is unique this is not possible, so the extremal function must have degree $d$.
		
		Next, let $P(z)=a_0+a_1 z + \dots + a_d z^d$ be a solution of \eqref{eq:extprob}. If $a_0=0$ then the polynomial $P(z)/z$ is also a solution of \eqref{eq:extprob}, and so we must have $a_0 \neq 0$. Define the polynomial $Q$ by $Q(z)=z^d \overline{P(1/\overline{z})}=\overline{a_d}+ \overline{a_{d-1}}z + \dots + \overline{a_0}z^d$. Then $Q$ is of degree $d$ with $Q(1)=P(1)=1$ and $\|Q\|_p=\|P\|_p$. Since the extremal function is unique it follows that $Q=P$. In particular if $re^{i t}$ (with $r>0$) is a zero of $P$, then so is $(1/r) e^{i t}$. 
		
		We now show that the zeroes of the extremal function must all be on the unit circle. Assume to the contrary that $re^{i t}$ is a zero of the extremal function $P$ and that $r \neq 1$, so that $(1/r) e^{i t}$ also is a zero of $P$. We will construct a polynomial $R$ of degree $d$ to show that $P$ cannot be extremal. To do so we let \[A(\theta)=(e^{i \theta}-re^{it})(e^{i \theta}-(1/r)e^{it}),\]
		and \[B(\theta)=(e^{i \theta}-1)^2.\]
		Assume first that $0<t<\pi$. Pick $\varepsilon>0$ so small that 
		\[|A(\theta)| > |A(\theta)-i\varepsilon B(\theta)|>0, \]
		for any $0<\theta<2 \pi$. To see that such an epsilon exists we first let \[m=\min_{\theta \in [0, 2\pi]} |A(e^{i\theta})|\] and observe that $m>0$ since $r \neq 1$. Using the reverse triangle inequality we see that for any $0<\varepsilon<m/4$ it follows that $|A(\theta)-i\varepsilon B(\theta)|>0$. Further, to show that 
		\begin{equation}\label{eq:eps2}
			|A(\theta)| > |A(\theta)-i\varepsilon B(\theta)|,
		\end{equation}
		we see that it suffices to show
		\[i(\overline{A(\theta)}B(\theta)-A(\theta)\overline{B(\theta)})-\varepsilon|B(\theta)|^2>0.\]  
		We see that
		\[\begin{split}
			&i(\overline{A(\theta)}B(\theta)-A(\theta)\overline{B(\theta)})-\varepsilon|B(\theta)|^2 \\ =&\frac{8 \sin(t)((1-r)^2+2r(1-\cos(\theta-t))) \sin^2(\theta/2)}{r}-\varepsilon16 \sin^4(\theta/2).
		\end{split}\]
		In particular, if we choose $\varepsilon<\sin(t)(1-r)^2/(2r)$ then \eqref{eq:eps2} holds. We now define the polynomial 
		\[R(z)=\frac{(z-re^{it})(z-(1/r)e^{it})-i\varepsilon (z-1)^2}{(z-re^{it})(z-(1/r)e^{it})} P(z),\] of degree $d$. Then $|R(z)|<|P(z)|$ for all $z \neq 1$ on the unit circle and $\|R\|_p < \|P\|_p$. It follows that $P$ cannot be an extremal function. This means that we have shown that if $re^{it}$ is a zero of the extremal function and $0<t<\pi$ then $r=1$. The arguments for the cases $t=\pi$ and $\pi<t<2\pi$ are similar. In the case $\pi<t<2\pi$ we can pick $\varepsilon>0$ so small that \[|A(\theta)| > |A(\theta)+i\varepsilon B(\theta)|>0, \]
		for any $0<\theta<2 \pi$ and adjust the definition of $R$ accordingly to show that $P$ cannot be extremal. For $t=\pi$, in fact for any $\pi/2 < t < 3\pi/2$, we can pick 
		$\varepsilon>0$ so small that \[|A(\theta)| > |A(\theta)+\varepsilon B(\theta)|>0, \]
		and again adjust the definition of $R$ accordingly to show that $P$ cannot be extremal. This means that all zeroes of the extremal function must have absolute value $1$. 
		
		Finally by considering the polynomial $S(z)=\overline{P(\overline{z})}$ it follows that if $P$ is extremal then so is $S$, and hence $P=S$. In particular, if $(z-e^{it_j})$ is a factor of $P$ then so is $(z-e^{-it_j})$. This implies that either $(z-e^{it_j})(z-e^{-it_j})$ is a factor of $P$ or $(z-e^{it_j})=(z-e^{-it_j})=(z+1)$ is  a factor of $P$. In particular, if $d$ is odd, $(z+1)$ must be a factor of $P$. The claimed structure follows.
	\end{proof}

\end{document}
